%=============================================================
%  TEMPLATE PRESENTATION BEAMER - KHOA HTTT - UIT
%  Dành cho sinh viên báo cáo môn học & khóa luận
%=============================================================
\documentclass[10pt,aspectratio=169]{beamer}
\usepackage[utf8]{vietnam}
\usepackage{amsmath,amssymb,amsfonts}
\usepackage{booktabs,array,longtable}
\usepackage{graphicx}
\usepackage{tikz}
\usetikzlibrary{shapes,arrows,arrows.meta,positioning,calc,
                shapes.geometric,fit,decorations.pathreplacing,
                decorations.markings,backgrounds,shadows}
\usepackage{pgfplots}
\pgfplotsset{compat=1.18}
\usepackage[many]{tcolorbox}
\usepackage{fontawesome5}
\usepackage{hyperref}
\usepackage{listings}

%=============================================================
%  1. TEMPLATE CUSTOMIZATION (Màu sắc, Logo và Thông tin)
%  Sinh viên thay đổi các thông số dưới đây để tùy biến slide
%=============================================================
\definecolor{uitblue}{RGB}{108,79,59}       % Màu chủ đạo (Nâu đậm của header/bìa)
\definecolor{uitdark}{RGB}{74,53,37}        % Màu phụ (Nâu tối ở dải bìa/footer)
\definecolor{accent}{RGB}{197,140,50}       % Màu nhấn mạnh (Vàng Gold)
\definecolor{accentlt}{RGB}{222,178,110}    % Màu nhấn phụ (Vàng cát)
\definecolor{greenteal}{RGB}{100,120,80}    % Màu ví dụ/thực hành (Sage Green)
\definecolor{lightgray}{RGB}{248,245,240}   % Màu nền ấm của slide (Cream/Ivory)
\definecolor{textgray}{RGB}{60,55,50}       % Màu chữ thân (Charcoal)

\newcommand{\titlebgimage}{title_bg.png}
\newcommand{\sectionbgimage}{section_bg.png} % Tên tệp ảnh nền trang ngăn cách section (nếu có)     % Tên tệp ảnh nền mờ trang bìa (nếu có)
\newcommand{\logonhomNC}{logo_nhom_nc.png}   % Logo nhóm nghiên cứu/Lab (nếu có)
\newcommand{\logonhomthuchien}{logo_nhom_thuchien.png} % Logo nhóm thực hiện/CLB (nếu có)
\newcommand{\gvhd}{}         % Giáo viên hướng dẫn (nếu có)
\newcommand{\slidewatermark}{logo_nhom_nc.png} % Tệp ảnh mờ làm hình nền từng slide (để trống nếu không dùng)
\newcommand{\contactemail}{ftisu@uit.edu.vn} % Email liên hệ hiển thị ở trang bìa

%=============================================================
%  2. THÔNG TIN BÁO CÁO (Metadata)
%=============================================================
\title[Distributed DL cho Calories]{Báo Cáo Thực Nghiệm: Distributed Deep Learning Cho Ước Lượng Calories}
\subtitle{Nutrition5k: RGB, RGB-D và Visual--Ingredient Fusion}
\author{Sinh viên thực hiện: An Văn Kết}
\date{Ngày báo cáo: \today}

%=============================================================
%  3. BEAMER THEME OVERRIDE (Cấu hình hiển thị - KHÔNG NÊN SỬA)
%=============================================================
\usetheme{default}
\usecolortheme{default}
\setbeamertemplate{navigation symbols}{}

% Thiết lập màu sắc các thành phần Beamer
\setbeamercolor{palette primary}{bg=uitblue,fg=white}
\setbeamercolor{frametitle}{bg=uitblue,fg=white}
\setbeamercolor{title}{bg=uitblue,fg=white}
\setbeamercolor{author in head/foot}{bg=uitdark,fg=white}
\setbeamercolor{title in head/foot}{bg=uitblue,fg=white}
\setbeamercolor{date in head/foot}{bg=uitdark!80,fg=white}
\setbeamercolor{section in head/foot}{bg=uitblue!70,fg=white}

\setbeamercolor{block title}{bg=uitblue!18,fg=uitdark}
\setbeamercolor{block body}{bg=uitblue!5,fg=textgray}
\setbeamercolor{alertblock title}{bg=accent!20,fg=accent!80!black}
\setbeamercolor{alertblock body}{bg=accent!5,fg=textgray}
\setbeamercolor{exampleblock title}{bg=greenteal!20,fg=greenteal!70!black}
\setbeamercolor{exampleblock body}{bg=greenteal!5,fg=textgray}
\setbeamercolor{item}{fg=uitblue}
\setbeamercolor{subitem}{fg=accent}
\setbeamercolor{enumerate item}{fg=uitblue}
\setbeamercolor{normal text}{fg=textgray}
\setbeamercolor{structure}{fg=uitblue}

% Định dạng tiêu đề slide (Frametitle)
% Thiết lập hình nền mờ dưới từng slide (Watermark)
\usebackgroundtemplate{
  \ifx\slidewatermark\empty\else
    \expandafter\IfFileExists\expandafter{\slidewatermark}{
      \begin{tikzpicture}[remember picture, overlay]
        \node[opacity=0.03, anchor=center] at (current page.center) {
          \includegraphics[width=0.3\paperwidth]{\slidewatermark}
        };
      \end{tikzpicture}
    }{}
  \fi
}

\setbeamertemplate{frametitle}{
  \vskip2pt
  \begin{beamercolorbox}[wd=\paperwidth,ht=0.56cm,dp=0.2cm,leftskip=0.4cm]{frametitle}
    \usebeamerfont{frametitle}\insertframetitle
    \ifx\insertframesubtitle\empty\else
      \;{\small\textcolor{accentlt}{$\triangleright$}}\;\usebeamerfont{framesubtitle}\textcolor{accentlt}{\insertframesubtitle}
    \fi
  \end{beamercolorbox}
  \vskip-2pt
  \begin{beamercolorbox}[wd=\paperwidth,ht=1.2pt,dp=0pt]{palette primary}
    \color{accent}\rule{\paperwidth}{1.2pt}
  \end{beamercolorbox}
}

% Định dạng chân trang (Footline)
\setbeamertemplate{footline}{
  \begin{beamercolorbox}[wd=\paperwidth,ht=0.4cm,dp=0.12cm]{author in head/foot}
    \hspace{0.3cm}{\scriptsize\faUniversity\; UIT -- Khoa Hệ thống Thông tin}
    \hfill
    {\scriptsize Báo cáo khoa học}
    \hfill
    {\scriptsize\insertframenumber/\inserttotalframenumber\hspace{0.3cm}}
  \end{beamercolorbox}
}

% Định dạng đầu trang (Headline)
\setbeamertemplate{headline}{
  \begin{beamercolorbox}[wd=\paperwidth,ht=0.06cm]{section in head/foot}
    \color{accent}\rule{\paperwidth}{0.06cm}
  \end{beamercolorbox}
}

% Định dạng ký hiệu danh sách (Itemize bullets)
\setbeamertemplate{itemize item}{\color{uitblue}\faCircle[regular]}
\setbeamertemplate{itemize subitem}{\color{accent}$\triangleright$}
\setbeamertemplate{enumerate item}{\color{uitblue}\insertenumlabel.}

% Định dạng trang bìa (Title Page)
\setbeamertemplate{title page}{
  \begin{tikzpicture}[remember picture, overlay]
    \fill[uitblue] (current page.north west) rectangle ([yshift=-0.58\paperheight]current page.north east);
    \fill[uitdark] ([yshift=-0.58\paperheight]current page.north west)
         rectangle ([yshift=-0.62\paperheight]current page.north east);
    \fill[lightgray] ([yshift=-0.62\paperheight]current page.north west)
         rectangle (current page.south east);
         
    \expandafter\IfFileExists\expandafter{\titlebgimage}{
      \node[opacity=0.10, anchor=center] at (current page.center) {\includegraphics[width=\paperwidth,height=\paperheight]{\titlebgimage}};
    }{}
         
    \node[white, anchor=north west, font=\large\bfseries, align=left, text width=0.9\paperwidth] (univ) at
         ([xshift=0.5cm, yshift=-0.4cm]current page.north west)
         {\faUniversity\quad TRƯỜNG ĐẠI HỌC CÔNG NGHỆ THÔNG TIN -- UIT \\
          \small\normalfont Khoa Hệ thống Thông tin $|$ Nhóm nghiên cứu FTISU};
          
    \node[white, anchor=north west, font=\LARGE\bfseries, text width=0.88\paperwidth, align=left] (title) at
         ([yshift=-0.25cm]univ.south west)
         {\inserttitle};
         
    \node[accentlt, anchor=north west, font=\normalsize\itshape, text width=0.88\paperwidth, align=left] (subtitle) at
         ([yshift=-0.15cm]title.south west)
         {\insertsubtitle};
         
    \node[uitdark, anchor=north west, font=\small, align=left] at
         ([xshift=0.8cm, yshift=-0.66\paperheight]current page.north west)
         {\faUser\quad \insertauthor \\
          \ifx\gvhd\empty\else
            \faGraduationCap\quad GVHD: \gvhd \\
          \fi
          \faEnvelope\quad \contactemail \\
          \faCalendar\quad \insertdate};
         
    % Trang trí dòng kẻ đứng
    \draw[accent,line width=4pt]
         ([xshift=0.5cm,yshift=-0.65\paperheight]current page.north west)
      -- ([xshift=0.5cm,yshift=-0.85\paperheight]current page.north west);

    % Hiển thị Logo nhóm nghiên cứu / câu lạc bộ / khoa ở góc dưới bên phải
        % Hiển thị Logo: Nhóm NC nằm bên TRÁI nhóm thực hiện (Góc dưới bên phải)
    \expandafter\IfFileExists\expandafter{\logonhomthuchien}{
        % Vẽ Logo nhóm thực hiện trước, neo vào góc và đặt tên node là 'logo2'
        \node[anchor=south east, inner sep=0pt] (logo2) at ([xshift=-0.5cm, yshift=0.4cm]current page.south east) {
            \includegraphics[height=0.7cm]{\logonhomthuchien}
        };
        
        \expandafter\IfFileExists\expandafter{\logonhomNC}{
            % Logo nhóm NC sẽ bám vào phía bên TRÁI (west) của logo2
            \node[anchor=east, inner sep=0pt] at ([xshift=-0.2cm]logo2.west) {
                \includegraphics[height=0.7cm]{\logonhomNC}
            };
        }{}
    }{
        % Trường hợp KHÔNG có logo nhóm thực hiện, chỉ có logo nhóm NC
        \expandafter\IfFileExists\expandafter{\logonhomNC}{
            \node[anchor=south east, inner sep=0pt] at ([xshift=-0.5cm, yshift=0.4cm]current page.south east) {
                \includegraphics[height=0.7cm]{\logonhomNC}
            };
        }{}
    }
    \end{tikzpicture}
}

\AtBeginSection[]{
  {
    \usebackgroundtemplate{
      \begin{tikzpicture}[remember picture, overlay]
        \expandafter\IfFileExists\expandafter{\sectionbgimage}{
          \node[anchor=center, inner sep=0pt] at (current page.center) {
            \includegraphics[width=\paperwidth,height=\paperheight]{\sectionbgimage}
          };
        }{}
        \fill[uitblue, opacity=0.92] (current page.north west) rectangle (current page.south east);
      \end{tikzpicture}
    }
    \begin{frame}[plain]
      \centering
      \vspace{1.8cm}
      {\color{white}\Huge\bfseries\thesection}
      \par\vskip 0.3cm
      {\color{accent}\rule{0.5\paperwidth}{2pt}}
      \par\vskip 0.4cm
      {\color{accentlt}\Large\bfseries\insertsectionhead}
    \end{frame}
  }
}

%=============================================================
%  4. CÁC HỘP THÔNG TIN TCOLORBOX
%=============================================================
\tcbuselibrary{skins,breakable}

% Hộp UIT mặc định (Màu xanh chủ đạo/nâu đất)
\newtcolorbox{uitbox}[2][]{
  enhanced,
  colback=uitblue!6, colframe=uitblue, coltitle=white,
  fonttitle=\bfseries\small, title={#2},
  boxrule=0.5pt, arc=4pt, left=4pt, right=4pt, top=3pt, bottom=3pt,
  attach boxed title to top left={yshift=-2mm,xshift=4mm},
  boxed title style={colback=uitblue, arc=3pt, boxrule=0pt},
  #1
}

% Hộp Accent nhấn mạnh (Màu vàng nhấn)
\newtcolorbox{accentbox}[2][]{
  enhanced,
  colback=accent!6, colframe=accent, coltitle=white,
  fonttitle=\bfseries\small, title={#2},
  boxrule=0.5pt, arc=4pt, left=4pt, right=4pt, top=3pt, bottom=3pt,
  attach boxed title to top left={yshift=-2mm,xshift=4mm},
  boxed title style={colback=accent, arc=3pt, boxrule=0pt},
  #1
}

% Hộp ví dụ (Màu xanh lá)
\newtcolorbox{greenbox}[2][]{
  enhanced,
  colback=greenteal!6, colframe=greenteal, coltitle=white,
  fonttitle=\bfseries\small, title={#2},
  boxrule=0.5pt, arc=4pt, left=4pt, right=4pt, top=3pt, bottom=3pt,
  attach boxed title to top left={yshift=-2mm,xshift=4mm},
  boxed title style={colback=greenteal, arc=3pt, boxrule=0pt},
  #1
}

% Hộp cảnh báo / hạn chế
\newtcolorbox{redbox}[2][]{
  enhanced,
  colback=red!4, colframe=red!60!black, coltitle=white,
  fonttitle=\bfseries\small, title={#2},
  boxrule=0.5pt, arc=4pt, left=4pt, right=4pt, top=3pt, bottom=3pt,
  attach boxed title to top left={yshift=-2mm,xshift=4mm},
  boxed title style={colback=red!60!black, arc=3pt, boxrule=0pt},
  #1
}

%=============================================================
%  5. CÁC MACRO RÚT GỌN (Tô màu từ khóa)
%=============================================================
\newcommand{\highlight}[1]{\textcolor{accent}{\bfseries #1}}
\newcommand{\kw}[1]{\textcolor{uitblue}{\bfseries #1}}

%=============================================================
%  PHẦN THÂN TÀI LIỆU
%=============================================================

\begin{document}

%--- Slide 1: Trang bìa ---
{
\usebackgroundtemplate{}
\begin{frame}[plain]
\vspace{0.25cm}
\begin{center}
\begin{beamercolorbox}[wd=0.94\paperwidth,ht=1.25cm,dp=0.28cm,center,rounded=true]{frametitle}
{\Large\bfseries Báo Cáo Thực Nghiệm: Distributed Deep Learning Cho Ước Lượng Calories}
\end{beamercolorbox}

\vspace{0.25cm}
{\large\color{accent}\bfseries Nutrition5k: RGB, RGB-D và Visual--Ingredient Fusion}

\vspace{0.55cm}
\begin{uitbox}[width=0.78\paperwidth]{Thông tin báo cáo}
\centering
\faUser\quad Sinh viên thực hiện: \textbf{An Văn Kết}\\[3pt]
\faUniversity\quad UIT -- Khoa Hệ thống Thông tin\\[3pt]
\faCalendar\quad Ngày báo cáo: \today
\end{uitbox}

\vspace{0.35cm}
\begin{accentbox}[width=0.78\paperwidth]{Mục tiêu chính}
\centering
Đánh giá \kw{accuracy--scalability trade-off} của các mô hình dự đoán calories khi huấn luyện 1 GPU và 2 GPU DDP.
\end{accentbox}
\end{center}
\end{frame}
}

%--- Slide 2 ---
\begin{frame}{Mục Tiêu Báo Cáo}
\scriptsize
\begin{columns}[T]
\column{0.33\textwidth}
\begin{uitbox}{Bài toán}
Dự đoán \kw{calories} của món ăn từ dữ liệu Nutrition5k.

\vspace{0.08cm}
Đây là bài toán \kw{regression}: đầu ra là một giá trị số kcal.
\end{uitbox}

\column{0.33\textwidth}
\begin{accentbox}{Vì sao cần Distributed?}
Huấn luyện mô hình ảnh và mô hình đa phương thức tốn nhiều thời gian, tính toán và bộ nhớ GPU.

\vspace{0.08cm}
Vì vậy cần kiểm tra khả năng mở rộng bằng \kw{Distributed Data Parallel}.
\end{accentbox}

\column{0.33\textwidth}
\begin{greenbox}{Câu hỏi chính}
Khi mô hình chuyển từ RGB đơn giản sang RGB-D hoặc Visual--Ingredient Fusion:

\vspace{0.05cm}
\begin{itemize}
  \item MAE có giảm không?
  \item 2 GPU có tăng tốc không?
  \item mô hình nào cân bằng tốt nhất?
\end{itemize}
\end{greenbox}
\end{columns}
\end{frame}

%--- Slide 3 ---
\begin{frame}{Dữ Liệu Và Thiết Lập Thí Nghiệm}
\scriptsize
\begin{columns}[T]
\column{0.48\textwidth}
\begin{uitbox}{Dữ liệu sử dụng}
\begin{itemize}
  \item Dataset: \kw{Nutrition5k}.
  \item Số mẫu có ảnh RGB hợp lệ: \highlight{3,244} món.
  \item Chia dữ liệu cố định: train/validation/test.
  \item Target: tổng \kw{calories} của mỗi món ăn.
\end{itemize}
\end{uitbox}

\vspace{0.12cm}
\begin{greenbox}{Thống kê calories}
\centering
\renewcommand{\arraystretch}{1.12}
\begin{tabular}{lrrr}
\toprule
\textbf{Split} & \textbf{Số mẫu} & \textbf{Mean} & \textbf{Std} \\
\midrule
Train & 2270 & 253.31 & 221.40 \\
Val & 487 & 256.67 & 212.09 \\
Test & 487 & 266.35 & 218.93 \\
\bottomrule
\end{tabular}
\end{greenbox}

\column{0.48\textwidth}
\begin{accentbox}{Thiết lập chung}
\begin{itemize}
  \item Epoch: \highlight{20}.
  \item Global batch size: \highlight{64}.
  \item 1 GPU: batch 64.
  \item 2 GPU: batch 32/GPU.
  \item Loss: \kw{SmoothL1Loss}.
  \item Optimizer: \kw{Adam}, learning rate $10^{-4}$.
  \item Thời gian đo: chỉ tính \kw{training time}, không tính validation.
\end{itemize}
\end{accentbox}
\end{columns}
\end{frame}

%--- Slide 4 ---
\begin{frame}{Vì Sao Bài Toán Phù Hợp Với Distributed Learning?}
\scriptsize
\begin{columns}[T]
\column{0.46\textwidth}
\begin{uitbox}{Cách phân tán đang dùng}
Hướng triển khai là \kw{Synchronous Data Parallelism} bằng PyTorch \kw{DistributedDataParallel}.

\vspace{0.08cm}
Mỗi GPU giữ một bản sao đầy đủ của mô hình, xử lý một phần mini-batch, sau đó đồng bộ gradient trước khi cập nhật tham số.
\end{uitbox}

\vspace{0.12cm}
\begin{greenbox}{Vì sao phù hợp?}
Ảnh món ăn là các sample độc lập, nên có thể chia batch cho nhiều GPU mà không làm thay đổi bản chất bài toán.
\end{greenbox}

\column{0.50\textwidth}
\begin{accentbox}{Luồng xử lý 2 GPU}
\centering
\begin{tabular}{c}
Dataset Nutrition5k \\
$\downarrow$ \\
\begin{tabular}{cc}
GPU 0 & GPU 1 \\
Model copy & Model copy \\
Batch 32 & Batch 32 \\
Local gradient & Local gradient \\
\end{tabular} \\
$\downarrow$ \\
Gradient synchronization / AllReduce \\
$\downarrow$ \\
Update cùng tham số mô hình
\end{tabular}
\end{accentbox}
\vspace{0.1cm}
\begin{redbox}{Điểm cần đánh giá}
Distributed không chỉ cần nhanh hơn, mà còn phải giữ được MAE/RMSE tương đương và sử dụng GPU hiệu quả.
\end{redbox}
\end{columns}
\end{frame}

%--- Slide 5 ---
\begin{frame}{Exp 1--3: RGB Baseline Và Kiểm Tra DDP}
\scriptsize
\begin{columns}[T]
\column{0.47\textwidth}
\begin{uitbox}{Ý nghĩa thí nghiệm}
\begin{itemize}
  \item \kw{Exp 1}: ResNet18 5 epoch, kiểm tra nhanh DDP có chạy được.
  \item \kw{Exp 2}: ResNet18 20 epoch, benchmark chính thức đầu tiên.
  \item \kw{Exp 3}: ResNet50 20 epoch, kiểm tra model RGB mạnh hơn.
\end{itemize}
\end{uitbox}

\vspace{0.12cm}
\begin{greenbox}{Kết luận nhanh}
ResNet50 cải thiện MAE rất mạnh so với ResNet18, nhưng thời gian huấn luyện cũng tăng.
\end{greenbox}

\column{0.50\textwidth}
\begin{accentbox}{Kết quả chính}
\centering
\resizebox{\linewidth}{!}{%
\begin{tabular}{lrrrr}
\toprule
\textbf{Model} & \textbf{GPU} & \textbf{Time} & \textbf{Throughput} & \textbf{Best MAE} \\
\midrule
ResNet18 & 1 & 443.32s & 102.41 & 209.96 \\
ResNet18 & 2 & 317.75s & 142.88 & 204.45 \\
ResNet50 & 1 & 550.37s & 82.49 & 77.41 \\
ResNet50 & 2 & 390.44s & 116.28 & 84.53 \\
\bottomrule
\end{tabular}}
\end{accentbox}
\vspace{0.1cm}
\begin{redbox}{Diễn giải}
2 GPU giúp speedup khoảng \highlight{1.40--1.41x}. Đây là bằng chứng bài toán có thể huấn luyện theo hướng distributed.
\end{redbox}
\end{columns}
\end{frame}

%--- Slide 6 ---
\begin{frame}{Exp 4: RGB-D Fusion}
\scriptsize
\begin{columns}[T]
\column{0.47\textwidth}
\begin{uitbox}{Mục tiêu}
Thêm thông tin \kw{depth} để mô hình hiểu tốt hơn về hình dạng và độ nổi của món ăn.

\vspace{0.08cm}
Depth được chuẩn hóa bằng percentile từ train set để giảm ảnh hưởng của outlier.
\end{uitbox}

\vspace{0.12cm}
\begin{accentbox}{Kiến trúc}
\centering
RGB $\rightarrow$ ResNet50 $\rightarrow$ 2048-d\\
Depth $\rightarrow$ ResNet18 $\rightarrow$ 512-d\\
$\downarrow$\\
Concat $\rightarrow$ Fusion FC $\rightarrow$ Calories
\end{accentbox}

\column{0.50\textwidth}
\begin{greenbox}{Kết quả}
\centering
\resizebox{\linewidth}{!}{%
\begin{tabular}{lrr}
\toprule
\textbf{Metric} & \textbf{1 GPU} & \textbf{2 GPU DDP} \\
\midrule
Train time & 800.92s & 632.78s \\
Throughput & 56.68 & 71.75 \\
Best Val MAE & 44.76 & \textbf{44.42} \\
Best Val RMSE & 73.59 & \textbf{69.94} \\
VRAM/GPU & 7.36GB & 3.86GB \\
\bottomrule
\end{tabular}}
\end{greenbox}
\vspace{0.1cm}
\begin{redbox}{Nhận xét}
Accuracy tốt nhất, nhưng speedup chỉ \highlight{1.27x}. Nguyên nhân có thể do thêm I/O depth, preprocessing và mô hình lớn hơn.
\end{redbox}
\end{columns}
\end{frame}

%--- Slide 7 ---
\begin{frame}{Exp 5: Visual--Ingredient Fusion}
\scriptsize
\begin{columns}[T]
\column{0.47\textwidth}
\begin{uitbox}{Mục tiêu}
Thêm danh sách nguyên liệu dưới dạng \kw{ingredient ID}, kết hợp với ảnh RGB để dự đoán calories.

\vspace{0.08cm}
Chỉ dùng \kw{ingr\_id}; không dùng grams, calories, fat, carb, protein để tránh data leakage.
\end{uitbox}

\vspace{0.12cm}
\begin{accentbox}{Kiến trúc}
\centering
RGB $\rightarrow$ ResNet50 $\rightarrow$ Visual feature\\
Ingredient IDs $\rightarrow$ Embedding $\rightarrow$ Mean pooling\\
$\downarrow$\\
Fusion FC $\rightarrow$ Calories
\end{accentbox}

\column{0.50\textwidth}
\begin{greenbox}{Kết quả}
\centering
\resizebox{\linewidth}{!}{%
\begin{tabular}{lrr}
\toprule
\textbf{Metric} & \textbf{1 GPU} & \textbf{2 GPU DDP} \\
\midrule
Train time & 545.45s & \textbf{368.11s} \\
Throughput & 83.23 & \textbf{123.33} \\
Best Val MAE & \textbf{45.88} & 49.15 \\
Test MAE & \textbf{49.85} & 52.30 \\
VRAM/GPU & 5.49GB & \textbf{3.10GB} \\
\bottomrule
\end{tabular}}
\end{greenbox}
\vspace{0.1cm}
\begin{redbox}{Nhận xét}
Accuracy gần RGB-D nhưng train nhanh hơn nhiều. Đây là mô hình cân bằng tốt nhất giữa accuracy và scalability.
\end{redbox}
\end{columns}
\end{frame}

%--- Slide 8 ---
\begin{frame}{So Sánh Scalability Giữa Các Kiến Trúc}
\scriptsize
\begin{accentbox}{Bảng khả năng mở rộng 1 GPU $\rightarrow$ 2 GPU}
\centering
\renewcommand{\arraystretch}{1.15}
\resizebox{0.96\linewidth}{!}{%
\begin{tabular}{lrrrrr}
\toprule
\textbf{Architecture} & \textbf{1 GPU} & \textbf{2 GPU} & \textbf{Speedup} & \textbf{Efficiency} & \textbf{Time Reduction} \\
\midrule
ResNet18 RGB & 443.32s & 317.75s & 1.40x & 69.8\% & 28.3\% \\
ResNet50 RGB & 550.37s & 390.44s & 1.41x & 70.5\% & 29.1\% \\
RGB-D Fusion & 800.92s & 632.78s & 1.27x & 63.3\% & 21.0\% \\
Visual--Ingredient & 545.45s & 368.11s & \textbf{1.48x} & \textbf{74.1\%} & \textbf{32.5\%} \\
\bottomrule
\end{tabular}}
\end{accentbox}

\vspace{0.12cm}
\begin{columns}[T]
\column{0.48\textwidth}
\begin{uitbox}{Điểm rút ra}
Mô hình nhiều thông tin hơn không tự động scale tốt hơn. RGB-D chính xác cao nhưng efficiency thấp hơn do chi phí I/O và xử lý depth.
\end{uitbox}
\column{0.48\textwidth}
\begin{greenbox}{Điểm tốt nhất}
Visual--Ingredient đạt speedup cao nhất và vẫn giữ MAE ở mức thấp trên test set.
\end{greenbox}
\end{columns}
\end{frame}

%--- Slide 9 ---
\begin{frame}{Kết Luận Và Hướng Đề Tài}
\scriptsize
\begin{columns}[T]
\column{0.48\textwidth}
\begin{uitbox}{Kết luận thực nghiệm}
\begin{itemize}
  \item Bài toán calories từ ảnh phù hợp để chạy \kw{Distributed Data Parallel}.
  \item 2 GPU giúp giảm thời gian huấn luyện khoảng \highlight{21--33\%}.
  \item Multimodal giúp giảm MAE mạnh so với RGB-only.
  \item RGB-D chính xác nhất trên validation; Visual--Ingredient cân bằng nhất về tốc độ và bộ nhớ.
\end{itemize}
\end{uitbox}

\column{0.48\textwidth}
\begin{accentbox}{Hướng nghiên cứu đề xuất}
\textbf{Accuracy--Scalability Trade-off} trong Distributed Multimodal Deep Learning cho ước lượng dinh dưỡng.

\vspace{0.10cm}
Các hướng có thể làm tiếp:
\begin{itemize}
  \item đo bottleneck I/O và GPU compute;
  \item thử gradient compression;
  \item đánh giá nhiều seed để tăng độ tin cậy;
  \item mở rộng từ calories sang calories + fat + carb + protein.
\end{itemize}
\end{accentbox}
\end{columns}
\end{frame}

\end{document}
